\documentclass[11pt,a4paper]{article}
\usepackage[T1]{fontenc}
\usepackage{lmodern,amsmath,amssymb,amsthm,booktabs,geometry,microtype}
\usepackage[hidelinks]{hyperref}
\geometry{margin=27mm}
\newtheorem{theorem}{Theorem}
\newtheorem{lemma}[theorem]{Lemma}
\newcommand{\R}{\mathrm R}
\newcommand{\U}{\mathrm U}
\newcommand{\Prob}{\mathbb P}
\title{An infinite family of fixed-cardinality counterexamples\\
to geodesic-probability monotonicity}
\author{Unsigned research draft}
\date{30 October 2026\\\small Public research draft; not peer reviewed}
\begin{document}
\maketitle
\begin{abstract}
For every integer $n\geq4$, we construct two up/right paths on the
$(n+1)\times(n+1)$ vertex grid, the second obtained from the first by a
boundary unwinding move. Exactly three nonendpoint vertices are selected
uniformly and assigned weight one; every other nonendpoint vertex has
weight zero. All tied maximizing paths count as geodesics. We prove by
elementary classification that the number of three-point configurations
making the unwound path optimal, minus the corresponding number for the
original path, is $-(5n^2-27n+32)/2$. The difference is strictly
negative for every $n\geq4$. Thus the fixed-cardinality analogue of a
path-probability ordering fails in arbitrarily large squares. This is
\emph{not} a counterexample to the independent Bernoulli all-parameter
conjecture: conditioning on the total number of occupied sites changes the
probability measure. The result extends a previously documented $n=4$
example to an analytically proved infinite family. Priority and independent
research significance remain unestablished.
\end{abstract}

\section{Setting and statement}
Let $V_n=\{0,\ldots,n\}^2$ and
$W_n=V_n\setminus\{(0,0),(n,n)\}$. Pick $S$ uniformly from the
three-element subsets of $W_n$, and assign weight $1$ to its elements
and weight $0$ to the other nonendpoint vertices. Since every up/right
path contains the two endpoints, arbitrary common endpoint weights do
not affect the maximizing set. A path $\gamma$ is a geodesic when its
score $|S\cap\gamma|$ is at least that of every up/right path; all ties
are counted.

Let $\R$ and $\U$ denote right and up steps. For $n\ge4$ define
\[
 \gamma_n=\R\U^2\R\U^{n-2}\R^{n-2},
 \qquad
 \pi_n=\U^2\R^2\U^{n-2}\R^{n-2}.
\]
Write $C_\eta(n)$ for the number of three-element sets $S\subset W_n$
for which $\eta$ is a geodesic.

\begin{theorem}\label{thm:main}
For every integer $n\ge4$, $\pi_n$ is obtained from $\gamma_n$ by a
genuine boundary unwinding move as in \cite[Section~2]{MN}, and
\[
 C_{\pi_n}(n)-C_{\gamma_n}(n)
 =-\frac{5n^2-27n+32}{2}<0.
\]
Consequently,
\[
 \Prob(\pi_n\text{ is a geodesic})-
 \Prob(\gamma_n\text{ is a geodesic})
 =-\frac{5n^2-27n+32}
 {2\binom{(n+1)^2-2}{3}}<0.
\]
\end{theorem}

In the signed-run convention of \cite{MN}, the source runs are
$(-1,2,-1,n-2,-(n-2))$. Their first two partial sums are $-1,1$.
Taking $l=0,k=2$ reverses the actual excursion prefix
$\R\U\U$ into $\U\U\R$; the next $\R$ merges with that prefix to give
$\pi_n$. The move is nonidentity and its interior excursion is nonempty.
In particular, we do not invoke the vacuous adjacent-corner case in
the printed definition.

The authors of \cite{MN} propose their ordering for independent
Bernoulli weights and prove strict monotonicity in a sufficiently small
positive parameter range. Our theorem concerns the \emph{conditional}
uniform three-point ensemble, which has dependent weights. It does
not compare independent Bernoulli probabilities. In particular,
a negative contribution from one cardinality layer need not determine
the sign after summing over all layers.

The case $n=4$ was already recorded in our earlier finite computation
\cite{B1}; it is not presented as a new example here. Our contribution
within this draft is the elementary analytic extension to the
specified infinite family, subject to the priority limitations in
Section~\ref{sec:novelty}.

\section{Reduction to three-point chain geometry}
Equip the grid with the coordinatewise partial order:
$(x,y)\leq(x',y')$ when $x\leq x'$ and $y\leq y'$.
A directed path can contain a set of vertices exactly when that set is
a chain in this order. Consequently, for a three-point set $S$, the
largest path score equals the maximum size of a chain contained in $S$.
This also proves that a path is optimal if and only if its own score
equals that chain length. No tie-breaking rule is used.

Introduce the four vertices where the two paths differ:
\[
 a=(1,0),\quad b=(1,1),\quad c=(0,1),\quad d=(0,2),
 \qquad
 L=\{a,b,c,d\}.
\]
The free vertices on $\gamma_n$ but not $\pi_n$ are $a,b$; those
on $\pi_n$ but not $\gamma_n$ are $c,d$. Let $C$ be their common
nonendpoint path vertices and $D=W_n\setminus(L\cup C)$. Explicitly,
\[
 C=\{(1,2)\}\cup\{(2,y):2\le y\le n\}
          \cup\{(x,n):3\le x\le n-1\}.
\]
All vertices in $C$ are above and to the right of all four vertices
in $L$. The sets $L,C,D$ partition $W_n$.

For each $S$ put
$\delta(S)=\mathbf 1_{\{\pi_n\text{ is optimal}\}}
             -\mathbf 1_{\{\gamma_n\text{ is optimal}\}}$.
The desired count difference is $\sum_{|S|=3}\delta(S)$.
Classify configurations according to $|S\cap L|$ and, when this equals
one, according to $|S\cap C|$. The resulting sums are
\[
\begin{array}{c|c|r}
 |S\cap L| & |S\cap C| & \sum\delta(S)\\ \hline
 0&\text{any}&0\\
 1&0&-\binom{n-1}{2}\\
 1&1&2-2(n-3)^2\\
 1&2&0\\
 2&\text{any}&1\\
 3&0&0
\end{array}
\tag{2.1}\label{eq:table}
\]
We derive each row below. Adding them gives
\[
 -\binom{n-1}{2}+2-2(n-3)^2+1
 =-\frac{5n^2-27n+32}{2}.
\]
For $n\ge4$, $5n^2-27n+32=(n-4)(5n-7)+4>0$.
This proves the numerical and sign assertions in Theorem~\ref{thm:main}.

\section{Proof of the configuration sums}
\subsection{Zero or three vertices in \texorpdfstring{$L$}{L}}
When $S\cap L=\varnothing$, both paths visit the same vertices of $S$,
hence their scores and success indicators agree.

If $S\subseteq L$, all four possible triples have longest chain length
two. The first two in the table are optimal only for $\gamma_n$, and
the last two only for $\pi_n$:
\[
\begin{array}{c|rrrr}
 S&\{a,b,c\}&\{a,b,d\}&\{a,c,d\}&\{b,c,d\}\\ \hline
 \delta(S)&-1&-1&1&1
\end{array}
\]
Their sum is zero.

\subsection{Exactly one vertex in \texorpdfstring{$L$}{L}}
Suppose first that $S=\{x,h_1,h_2\}$ with $x\in L$ and
$h_1,h_2\in C$. The two common vertices lie on both paths and the
local vertex lies on exactly one. That path collects all three ones
and is necessarily optimal; the other collects two and is not.
Since two choices of $x$ favor each path, the total difference is zero.

Next suppose $S=\{x,z_1,z_2\}$ with $x\in L$ and $z_1,z_2\in D$.
Only the path containing $x$ can be optimal, and it is optimal
exactly when $S$ is a three-point antichain. The counts of unordered
pairs $\{z_1,z_2\}\subset D$ satisfying this condition are
\[
 \begin{array}{c|rrrr}
 x&a&b&c&d\\ \hline
 \#&0&(n-2)(n-1)&0&\binom{n-1}{2}
 \end{array}.
\]
Indeed, points in $D$ incomparable with $a$ are in the single
column $(0,y)$, $3\le y\le n$; points incomparable with $c$ are
in the single row $(x,0)$, $2\le x\le n$. Each such set is itself
a chain. For $b$, the incomparable set is the union of that column
and row, and a pair must take one point from each. For $d$, the
incomparable points comprise rows $y=0,1$ with $2\le x\le n$;
an antichain pair takes one from each row with the lower-row point
strictly to the right. Thus the difference is
\[
 0+\binom{n-1}{2}-0-(n-2)(n-1)
 =-\binom{n-1}{2}.
\]

Finally let $S=\{x,h,z\}$ with $x\in L$, $h\in C$, $z\in D$.
The path containing $x$ collects two ones, the other collects one,
and the former is optimal exactly when the three vertices do not
form a chain. For each fixed $h$, the constant numbers of choices
$z\in D$ cancel between the two positive and two negative choices
of $x$. Hence the contribution is
\begin{equation}\label{eq:chains}
\sum_{h\in C}\sum_{z\in D}
 \bigl(1_{\{a,h,z\}\ \mathrm{chain}}
      +1_{\{b,h,z\}\ \mathrm{chain}}
      -1_{\{c,h,z\}\ \mathrm{chain}}
      -1_{\{d,h,z\}\ \mathrm{chain}}\bigr).
\end{equation}
For every $x\in L$ and $h\in C$, $x\le h$. The triple is a chain
exactly if $z\le x$, $x\le z\le h$, or $z\ge h$. No vertex of
$D$ lies below any $x\in L$, and the $z\ge h$ term is independent
of $x$ and cancels in \eqref{eq:chains}. Only the rectangles remain.

For $h=(h_x,h_y)\in C$, define
\[
 T(h)=\sum_{z\in D}
 \left(1_{a\le z\le h}+1_{b\le z\le h}
       -1_{c\le z\le h}-1_{d\le z\le h}\right).
\]
In the full grid, before deleting $L,C$ and the endpoints, the
signed rectangle count is
\[
 h_x(h_y+1)+h_xh_y-(h_x+1)h_y-(h_x+1)(h_y-1)
 =2h_x-2h_y+1.
\]
Each common-path vertex contributes zero to the signed sum;
the origin contributes zero, and the terminal endpoint is not in
any of these rectangles. The four local vertices have signed
contributions $1,1,-1,-2$, respectively, summing to $-1$.
Removing them therefore yields
\[
 T(h)=2(h_x-h_y+1).
\]
Summing over the explicit common-path set gives
\[
 \sum_{h\in C}T(h)
 =2-(n-3)(n-2)-(n-4)(n-3)
 =2-2(n-3)^2.
\]
This proves the remaining one-local-vertex row.

\subsection{Exactly two vertices in \texorpdfstring{$L$}{L}}
If one of the two vertices is in $\{a,b\}$ and the other is in
$\{c,d\}$, the paths have the same score. It remains to compare
$S=\{a,b,z\}$ with $S=\{c,d,z\}$, where $z\in C\cup D$.
For $z\in C$ the two configurations both contribute success to the
path containing the local pair, so their total contributions cancel.

For $z\in D$, the path containing the local pair collects two ones
and is optimal unless the triple is a chain. The unrestricted numbers
of choices of $z$ again cancel. Let $N_{ab}$ and $N_{cd}$ be the
numbers of $z\in D$ forming chains with $\{a,b\}$ and $\{c,d\}$.
All points of $D$ below these local pairs have been excluded. Thus
\[
 N_{ab}-N_{cd}
 =\#\{(x,1):2\le x\le n\}
  -\#\{(0,y):3\le y\le n\}
 =(n-1)-(n-2)=1.
\]
This is the last nonzero sum in \eqref{eq:table} and completes the proof.

\section{Independent computational checks}
The analytic argument above is finite symbolic counting, not an
extrapolation from examples. As a separate check, the supplied
\path{verify_sparse_dp.py} enumerates every three-point set for
$n=4,\ldots,11$. For each configuration it computes the longest
chain using an independently written rank dynamic program on the
three occupied points, and scores each path using bit masks.
It verifies the total difference and every row of \eqref{eq:table}.
It does not enumerate all $2^{(n+1)^2-2}$ environments.

For $n=5$ the saved \path{n5_independent_check.json} also compares
the complete set of $5{,}984$ three-point environments for all
$252$ paths, using separate lattice-dynamic-programming and
point-chain implementations. The two implementations agree on
every path count; two separately written move generators agree on
all $560$ move witnesses, of which $12$ have negative difference.
This is finite implementation evidence, not a proof-assistant
formalization or a substitute for independent human peer review.

\begin{center}
\begin{tabular}{rrrrr}
\toprule
$n$&$C_{\gamma_n}$&$C_{\pi_n}$&difference&three-point environments\\
\midrule
4&309&307&$-2$&1,771\\
5&864&853&$-11$&5,984\\
6&2,010&1,985&$-25$&16,215\\
7&4,122&4,078&$-44$&37,820\\
8&7,707&7,639&$-68$&79,079\\
\bottomrule
\end{tabular}
\end{center}
See \path{verify_sparse_dp.json} for the full $n=4,\ldots,11$ replay.
To reproduce the replay, run \texttt{python verify\_sparse\_dp.py} from
the \path{fixed_k_followup} directory. Python's standard library is
sufficient. Run without \texttt{-O}, as assertions are part of the checks.

\section{Prior work and limitations}\label{sec:novelty}
The prior geometric path-order questions and maximizing-path setting
are discussed by Alberts and Cator \cite[Section~8]{AC}. The specific
bump/unwind operations and small-parameter Bernoulli result are due to
McKeown and Nakajima \cite{MN}. The $n=4$ fixed-three-point
counterexample appears in the earlier finite companion \cite{B1}.

A bounded search for fixed-cardinality geodesic comparisons, the
explicit path words and the candidate polynomial did not locate the
same infinite-family formula. The searched sources and access limits
are recorded in \path{literature_audit_20261030.txt}. Forward citation
coverage and the global literature review are incomplete, so exact
priority is not established. Degree elevation, fixed-cardinality
decomposition and longest-chain counting are standard methods; no
claim of a new general technique is made.

The proof addresses only the displayed path family and exactly three
occupied nonendpoint sites. It does not establish the full Bernoulli
conjecture, classify other fixed-cardinality levels, or claim
publication-quality significance. This unsigned draft has not been
human peer reviewed or formally verified. It is deposited as a follow-up in the companion GitHub repository.
No journal or arXiv submission is represented by that deposit.

\begin{thebibliography}{9}
\bibitem{MN} S. McKeown and S. Nakajima,
\emph{On the most likely geodesic in last passage percolation},
arXiv:2608.11378v2, 20 September 2026,
\url{https://arxiv.org/html/2608.11378v2}.
\bibitem{AC} T. Alberts and E. Cator,
\emph{On the Passage Time Geometry of the Last Passage Percolation Problem},
ALEA \textbf{18} (2021), 211--247,
\url{https://alea.impa.br/articles/v18/18-10.pdf}.
\bibitem{B1} Unsigned finite research draft,
\emph{An exact all-parameter verification of the bump--unwind path order
in Bernoulli last-passage percolation for $n=4$},
2026, Section~4.1,
\url{https://github.com/SeraphimTsui/bernoulli-lpp-n4}.
\end{thebibliography}
\end{document}
